% ======================================================================
% Capítulo: Conclusiones y líneas futuras
% Extensión orientativa: 3--4 páginas.
% ======================================================================
\chapter{Conclusiones y líneas futuras}
\label{cap:conclusiones}

\section{Cumplimiento de objetivos}
\label{sec:cumplimiento}
% TODO: OE1--OE7 uno a uno, señalando cambios de alcance.

\section{Aportaciones}
\label{sec:aportaciones}
% TODO: Arena y multi-tenancy como base de validación; métricas de coste para elegir modelo.

\section{Lecciones aprendidas}
\label{sec:lecciones}
% TODO: Técnicas y de proceso (SDD con agentes de IA).

\section{Líneas futuras}
\label{sec:futuras}
% TODO: Boletines provinciales, alertas de subvenciones y plazos, observabilidad de nivel 2, apertura a la ciudadanía.

\paragraph{Clave de API propia en producción.} La plataforma permite que cada ayuntamiento configure su proveedor y su clave de API para el chat, mientras la ingesta y los \emph{embeddings} se mantienen centralizados (sección~\ref{sec:clave-propia}). Hoy no se usa por los costes; activarlo permitiría que cada entidad pague solo sus conversaciones. Quedan como mejoras posibles llevar las claves a un almacén de secretos externo y valorar si compensa pasar los \emph{embeddings} de la plataforma a un proveedor externo, lo que exigiría reindexar todos los boletines con el nuevo modelo.
